% GENERATED FILE. Edit canonical lessons, then run npm run generate:books.
% canonical-source-hash: fnv1a64:f43d1eb80f1b7a0d
% canonical-lessons: IT-C264-poi, IT-C264-raccontare, IT-C264-lettura, IT-C264-ripasso-racconto

\chapter{First, Then, In the End --- A Story Told in Order}
\label{ch:it-a-story-told-in-order}
\hlchaptermodality{264}

\begin{chapteropening}
\textbf{By the end of this chapter:} \emph{I can tell a short connected story in order with prima, poi and alla fine, and read one of about a hundred words.}

Complete IT-C264-ripasso-racconto: retell the reading's journey aloud, choosing scene or event for every verb.
\end{chapteropening}

\section[prima, poi, alla fine]{prima, poi, alla fine — the events, in order}

\label{lesson:IT-C264-poi}



\begin{quote}
Chapter twenty-eight put a dinner in order: \emph{prima} the first course, \emph{dopo} the second. Say \textquotedblleft{}before\textquotedblright{} and \textquotedblleft{}after\textquotedblright{}. (\emph{Prima}, \emph{dopo}.)

Now the same words go to work on a story. The scene is painted; the events need a running order.
\end{quote}

\begin{grammarlens}[title={Grammar Lens: three places in a story}]
\begin{quote}
\textbf{Prima} ho comprato il biglietto. — First I bought the ticket. \textbf{Poi} ho aspettato il treno. — Then I waited for the train. \textbf{Alla fine} sono tornato a casa. — In the end I went back home.
\end{quote}

\textbf{Poi} is the plain \textquotedblleft{}then, next\textquotedblright{}. It stands between two events and points from the first to the second, and you can use it as many times as the story has events: \emph{poi … poi … poi}.

\textbf{Dopo} does a little more. It can stand alone, \textquotedblleft{}afterwards\textquotedblright{}, but it can also take a noun: \emph{dopo il caffè}, \textquotedblleft{}after the coffee\textquotedblright{}. \emph{Poi} cannot.

\textbf{Alla fine} is \textquotedblleft{}in the end\textquotedblright{}: the last event, and usually the outcome. You have \textbf{infine}, \textquotedblleft{}finally\textquotedblright{}, too; \emph{infine} closes a list, \emph{alla fine} closes a story. \emph{C'era molto traffico, ma alla fine sono tornato a casa} — there was a lot of traffic, but in the end I got home.

Notice what these words do \textbf{not} touch: the scene. \emph{Faceva freddo} needs no \emph{poi}. Order belongs to the events.
\end{grammarlens}

\begin{cousinweb}
\textbf{poi} comes from Latin \emph{post}, \textquotedblleft{}after, behind\textquotedblright{}, through a Vulgar Latin \emph{pos} that had already dropped the \emph{t} (EtymDB 2.1). And you met \emph{post} before: \textbf{dopo} is \emph{dē} + \emph{post}, two Latin words welded together. So \emph{poi} and \emph{dopo} are the same root, once bare and once with \emph{dē} in front. English keeps \emph{post} in \textbf{postpone} and \textbf{postscript}.

\textbf{fine}, \textquotedblleft{}end\textquotedblright{}, in \emph{alla fine}, is Latin \emph{fīnis}, \textquotedblleft{}boundary, end\textquotedblright{} — inherited (EtymDB 2.1). English \textbf{final}, \textbf{finish} and \textbf{define} share it.
\end{cousinweb}

\subsection*{Your turn}
Say these aloud:

\begin{itemize}
\raggedright
  \item \emph{Prima ho comprato il biglietto. Poi ho aspettato il treno.}
  \item \emph{dopo il caffè} — then try \emph{poi} in its place, and hear why it fails
  \item \emph{Alla fine sono tornato a casa.}
  \item \emph{poi} and \emph{dopo}, then English \emph{postpone}
\end{itemize}

\begin{tcolorbox}[breakable,colback=teal!4,colframe=teal!35!black,title={Before you move on}]
Which of \emph{poi} and \emph{dopo} can take a noun? (\emph{\textbf{Dopo}} — \emph{dopo il caffè}.) Which closes a story, \emph{infine} or \emph{alla fine}? (\emph{\textbf{Alla fine.}}) Which Latin word is in both \emph{poi} and \emph{dopo}? (\emph{\textbf{Post}}, \textquotedblleft{}after\textquotedblright{}.)
\end{tcolorbox}
\section[raccontare]{raccontare — to tell a story, which is to count it out}

\label{lesson:IT-C264-raccontare}



\begin{quote}
Three pieces are in your hands: the imperfetto to hold the scene still, \emph{prima}, \emph{poi} and \emph{alla fine} to order the events, and \emph{perché} to explain one of them. This lesson adds no fourth piece. It puts the three together, under the verb that names the act.
\end{quote}

\begin{grammarlens}[title={Grammar Lens: the backbone of a told story}]
\textbf{raccontare}, \textquotedblleft{}to tell (a story)\textquotedblright{}, is a plain \emph{-are} verb you have had since chapter one hundred and thirty-six. Here is what it tells, once it has both pasts to work with:

\begin{quote}
\textbf{Era} sabato. \textbf{Faceva} caldo e \textbf{c'era} il sole. \textbf{Prima} \textbf{sono andato} al mercato. \textbf{Poi} \textbf{ho incontrato} un amico, e \textbf{abbiamo parlato} molto. \textbf{Alla fine} \textbf{sono tornato} a casa, \textbf{perché ero} stanco.
\end{quote}

\begin{quote}
\emph{It was Saturday. It was hot and the sun was out. First I went to the market. Then I met a friend, and we talked a lot. In the end I went home, because I was tired.}
\end{quote}

Follow the two pasts as they divide the work. \emph{Era}, \emph{faceva}, \emph{c'era}: imperfetto, the scene, held still. \emph{Sono andato}, \emph{ho incontrato}, \emph{abbiamo parlato}, \emph{sono tornato}: passato prossimo, the events dropped in. \emph{Prima}, \emph{poi} and \emph{alla fine} put the events in order; \emph{perché} explains the last one, and its reason, \emph{ero stanco}, is scene again. Every word of it is one this book has taught you.
\end{grammarlens}

\begin{cousinweb}
\textbf{raccontare} is built on \textbf{contare}, \textquotedblleft{}to count\textquotedblright{}, which you met in chapter two hundred and six. \emph{Contare} is inherited from Latin \emph{computāre}, \textquotedblleft{}to reckon up\textquotedblright{} (EtymDB 2.1), and \emph{computāre} is \emph{com-}, \textquotedblleft{}together\textquotedblright{}, plus \emph{putāre}, \textquotedblleft{}to settle a reckoning\textquotedblright{} (Lewis \& Short). The \emph{ra-} in front is Italian's \emph{re-} plus \emph{a-}: to count something out for someone. EtymDB 2.1 has no row for \emph{raccontare} itself, so that last step is the usual reading of its parts.

English has the same picture. To \textbf{recount} a story is to count it out (neither source traces that English word), and \textbf{tell} once meant \textquotedblleft{}to count\textquotedblright{} (EtymDB 2.1) — a bank \textbf{teller} still does.
\end{cousinweb}

\subsection*{Your turn}
Say these aloud:

\begin{itemize}
\raggedright
  \item the scene — \emph{Era sabato. Faceva caldo e c'era il sole.}
  \item the events — \emph{Prima sono andato al mercato. Poi ho incontrato un amico.}
  \item the end, with its reason — \emph{Alla fine sono tornato a casa, perché ero stanco.}
  \item \emph{raccontare}, \emph{contare}, then English \emph{recount}
\end{itemize}

\begin{tcolorbox}[breakable,colback=teal!4,colframe=teal!35!black,title={Before you move on}]
In your story, which past held the sunshine and which one took you to the market? (\emph{\textbf{C'era}} — imperfetto; \emph{\textbf{sono andato}} — passato prossimo.) Which words put the events in order? (\emph{\textbf{Prima}}, \emph{\textbf{poi}}, \emph{\textbf{alla fine}}.) What verb is inside \emph{raccontare}? (\emph{\textbf{Contare}}, \textquotedblleft{}to count\textquotedblright{}, from Latin \emph{computāre}.)
\end{tcolorbox}
\section[(reading)]{(reading) — a journey, told in both pasts}

\label{lesson:IT-C264-lettura}



\begin{quote}
Chapter thirty-six gave you forty-seven words in a row, all in the present. This passage is twice as long, and it happened yesterday.

Every word in it was taught before this page.
\end{quote}

\begin{tcolorbox}[breakable,skin=enhanced,colback=orange!6,colframe=orange!45!black,boxrule=0.5pt,arc=1mm,left=6pt,right=6pt,top=4pt,bottom=4pt,fonttitle=\bfseries,title={Read this}]
\begin{quote}
Ieri mattina faceva freddo e pioveva. Ero alla stazione, e il mio amico era lì. C'erano molte persone, e il treno era in ritardo. Prima abbiamo aspettato al binario. Poi abbiamo preso un caffè, perché eravamo stanchi. Mentre parlavamo, hanno annunciato il treno, e siamo saliti subito.
\end{quote}

\begin{quote}
Il viaggio era lungo, ma non era noioso: il mio amico raccontava la storia di suo padre, e io ascoltavo. Fuori pioveva ancora, ma in treno faceva caldo. Alla fine, la sera, eravamo al mare. Il mare era bello, e c'era la luna. Ero stanco, ma ero contento.
\end{quote}

Read it once without stopping. Do not translate as you go; let the sentences arrive.

Then read it a second time and sort every verb into one of two piles: the scene, or what happened.
\end{tcolorbox}

\begin{culture}
The events pile is short: \emph{abbiamo aspettato}, \emph{abbiamo preso}, \emph{hanno annunciato}, \emph{siamo saliti}. Four things happened, and \emph{prima} and \emph{poi} put the first two in order.

Everything else is scene: the cold, the rain, the crowd, the late train, the moon. A story is mostly background, and the imperfetto marks every piece of it.

\emph{Mentre parlavamo} is still running when the announcement cuts into it. And \emph{raccontava} is scene too: the telling went on while the train moved. \emph{Ha raccontato} would have made it one finished event.
\end{culture}

\subsection*{Your turn}
\begin{itemize}
\raggedright
  \item \emph{Say it:} the passage aloud, once, without stopping
\end{itemize}

\begin{tcolorbox}[breakable,colback=teal!4,colframe=teal!35!black,title={Before you move on}]
How many words in that passage were new? (\textbf{None.})

Why is \emph{raccontava} in the imperfetto? (The telling \textbf{went on} while the train moved: it is scene.)

Which word opened the line where a scene and an event met? (\emph{\textbf{Mentre.}})
\end{tcolorbox}
\section[(review)]{Review — tell it back}

\label{lesson:IT-C264-ripasso-racconto}



\begin{quote}
Say \textquotedblleft{}first\textquotedblright{}, \textquotedblleft{}then\textquotedblright{} and \textquotedblleft{}in the end\textquotedblright{}. (\emph{Prima, poi, alla fine.})
\end{quote}

\begin{grammarlens}[title={What you've built}]
You have read a story. Now you tell one, which is the same skill turned round: a reader sorts verbs into scene and event, and a teller has to choose them.

The journey from the reading, in five cues. Answer each one in Italian, out loud, in a full sentence. Choose the past as you go: is this cue the scene, or something that happened?
\end{grammarlens}

\subsection*{Your turn}
Say these aloud:

\begin{itemize}
\raggedright
  \item the weather that morning (\emph{Ieri mattina faceva freddo e pioveva.})
  \item where you were, and who was with you (\emph{Ero alla stazione, e il mio amico era lì.})
  \item the first thing you did, then the second (\emph{Prima abbiamo aspettato al binario. Poi abbiamo preso un caffè.})
  \item what cut into your talk (\emph{Mentre parlavamo, hanno annunciato il treno.})
  \item where you were in the end, and how you felt (\emph{Alla fine eravamo al mare. Ero stanco, ma ero contento.})
\end{itemize}

Now tell the whole journey once more, in your own order and your own words, without the cues.

\begin{tcolorbox}[breakable,colback=teal!4,colframe=teal!35!black,title={Before you move on}]
Which of your five cues were scene? (The weather, where you were, the talk that was running, how you felt.) Which were events? (Waiting, the coffee, the announcement.) Which Latin word is behind \emph{poi}? (\emph{\textbf{Post}}.) And what are you doing, literally, when you \emph{raccontare}? (\textbf{Counting it out} — \emph{contare}.)
\end{tcolorbox}
